\section{Introduction}
\label{sec:introduction}

The AI Act~\cite{ai_act} moves the providers of high-risk AI systems from declared to demonstrable compliance. Systems listed in its Annex III, such as creditworthiness assessment, recruitment and access to essential services, must meet requirements that can only be shown through records: technical documentation kept up to date (Art.~11 and Annex~IV), automatic logging of events over the system's lifetime (Art.~12), human oversight (Art.~14), post-market monitoring (Art.~72) and the reporting of serious incidents (Art.~73). Whoever supervises such a system will read those records and act on them.

The records that machine-learning pipelines keep were not made for that reader. MLOps tools log for operational observability: their records are mutable, held by the organisation that produces them, and not linked to any obligation. A human approval is a flag in a database, not a statement with a verifiable author. An external auditor who receives such records must trust the party under audit to have kept them intact, which is the old problem of the integrity of the witness.

The timing makes the problem practical. Regulation (EU) 2026/1744~\cite{omnibus_2026} set the application of the Annex III requirements to 2 December 2027 and of those under Art.~6(1) to 2 August 2028, and left the duties on logging, oversight, accuracy and incident reporting as they were. Implementing Regulation (EU) 2025/2531~\cite{qel_2025} fixed the reference standards for qualified electronic ledgers under eIDAS. The European Data Protection Board adopted the final version of its guidelines on blockchain in July 2026~\cite{edpb_2026}. And the new product liability directive~\cite{pld_2024}, which covers software, applies to products placed on the market after 9 December 2026, with a duty to disclose evidence in litigation.

The technical pieces exist, but apart. Tamper-evident logging is mature, from secure audit logs~\cite{schneier1999} and history trees~\cite{crosby2009} to Certificate Transparency~\cite{rfc9162}, yet it is not tied to legal obligations. Work aimed at the AI Act maps obligations to controls and documentation~\cite{lucaj2025,ugarte2026,buscemi2026,palumbo2025,ramos2024}, yet none of it signs or anchors the evidence when the pipeline produces it. Among the works we reviewed, none states what its evidence does not establish, and none measures where the cost of verifiable evidence lies.

This article joins the two lines. Following design science research with an artificial ex post evaluation, it presents a compliance-as-code architecture that turns AI Act obligations into evidence records a third party can verify, produced as a by-product of the pipeline. It answers three research questions:
\begin{description}[leftmargin=0pt, itemsep=2pt]
\item[RQ1.] How can AI Act obligations for high-risk systems be translated into evidence records that third parties can verify, and what does each record fail to establish?
\item[RQ2.] Which integrity, minimisation and selective-disclosure guarantees can such records offer, and which threats remain outside them?
\item[RQ3.] What does verifiable evidence cost, and where does the cost lie: in cryptography, in persistence or in distributed coordination?
\end{description}

The article makes five contributions.
\begin{enumerate}[label=C\arabic*., leftmargin=*, itemsep=2pt]
\item A derivation method from obligation to pipeline requirement, evidence and mechanism, applied to seven provisions in a matrix with an explicit column on what each record does not establish, and an evidence model that commits to the artefact, keeps the inputs of each decision and correlates the events of a run (Section~\ref{sec:derivation}).
\item A Compliance Policy Engine with fail-closed validation, a closed rule vocabulary, policies archived by hash and decisions that can be re-executed from the record, writing Ed25519-signed, hash-chained records (Section~\ref{sec:architecture}).
\item An OSCAL 1.1.2 profile of assessment results that separates technical objectives from legal provisions, with a redacted mode for Art.~78 that keeps the integrity anchors (Section~\ref{sec:architecture}).
\item An evaluation that locates the cost: a four-arm ablation, per-record and batched anchoring on Hyperledger Fabric, a credit-scoring pipeline on real data, an industrial data stream, a tamper experiment and fourteen audit questions (Section~\ref{sec:evaluation}).
\item Interoperability lessons from eight defects found between the Python issuer and the Go chaincode, and a shared boundary matrix that works as a conformance test (Section~\ref{sec:eval-interop}).
\end{enumerate}

The artefact extends the prototype of a master's dissertation~\cite{sousa2026}. The second version of the evidence model, its canonical form, signed approvals, batched anchoring and every experiment reported here are new to this article.

The scope is deliberately narrow. The evidence establishes the integrity and provenance of what was recorded; it does not establish that the system complies with the AI Act, that the recorded metrics were true, that no event was withheld, or that a court would give the records any particular weight. Section~\ref{sec:background} reviews related work, Section~\ref{sec:approach} sets out the research approach, Section~\ref{sec:derivation} derives the evidence from the obligations, Section~\ref{sec:architecture} describes the architecture and Section~\ref{sec:evaluation} evaluates it. Section~\ref{sec:discussion} discusses what the results mean, Section~\ref{sec:threats} the threats to their validity, and Section~\ref{sec:conclusion} concludes.
